%% kapa: exam=1 round=26 year=2015 period=1 form=A subject=civil_law subject_ko=민법 seq=35 number=35
%% source: Q-Net 감정평가사 기출문제 (cst00309) artlSeq=5138001 file=extracted/1008/5138001/제26회_감정평가사 1차_기출/제26회_감정평가사 1차_기출/1교시_감평_A형.pdf
%% conversion: kordoc 4.24.51 -> tools/gen1.py (auto); review=proofread-sonnet; flags=none
\begin{question}{35}
「민법」상 공유에 관한 설명으로 옳지 않은 것은? (특약은 고려하지 않고, 다툼이 있으면 판례에 따름)
\begin{choices}
  \item 각 공유자는 자기 지분을 자유롭게 처분할 수 있다.
  \item 소수지분권자가 다른 공유자의 동의 없이 공유물을 배타적으로 점유하는 경우, 다른 소수지분권자는 그 점유자를 상대로 보존행위에 기하여 공유물의 인도를 청구할 수 없다.
  \item 공유자 1인이 포기한 지분은 다른 공유자에게 각 지분의 비율로 귀속한다.
  \item 분할에 관한 협의가 성립한 후에 공유물분할소송을 제기하는 것은 허용되지 않는다.
  \item 공유자는 다른 공유자가 공유물 분할로 인하여 취득한 물건에 대하여 그 지분의 비율로 매도인과 동일한 담보책임이 있다.
\end{choices}
\end{question}
\answer{2}
